\documentclass[11pt]{article}
\usepackage{amsmath,amssymb,amsthm,hyperref}
\newtheorem{theorem}{Theorem}
\newtheorem{lemma}{Lemma}
\newcommand{\E}{\mathbb E}
\title{Arithmetic balancing using logarithmically many fair-block profiles}
\author{Argument by the rational-design research agent; written and checked independently}
\date{30 September 2026}
\begin{document}\maketitle
\noindent\textbf{Scope.} This note concerns only the arithmetic of block
multiplicities. It does not prove that concatenating the resulting
fair blocks preserves permutation fairness, and therefore does not
resolve polynomial equalization. Indeed, $n$ cyclic relabelings of any
word already balance its multiplicities, with total length $n$ times
the original length. The point here is to reduce the number of block
profiles to $O(\log n)$ while retaining polynomial total length.

\begin{lemma}\label{gcd}
Let $c_1,\ldots,c_n$ be positive integers and $q\ge1$. There are
permutations $\sigma_1,\ldots,\sigma_q$ for which
\[
 L:=\operatorname{lcm}_{i\in[n]}
       \gcd(c_{\sigma_1(i)},\ldots,c_{\sigma_q(i)})
 \le\left(\prod_{i=1}^n c_i\right)^{n^{1/q}/n}.
\]
\end{lemma}
\begin{proof}
Choose the $q$ permutations independently and uniformly.
For a prime $p$ and $a\ge1$, let $k_{p,a}$ count the indices $j$
with $p^a\mid c_j$, and put $x=k_{p,a}/n$.
The event $p^a\mid L$ requires at least one coordinate $i$ whose
$q$ selected multiplicities are all divisible by $p^a$.
The union bound gives
\[
 \Pr(p^a\mid L)\le\min(1,nx^q)\le n^{1/q}x.
\]
The final inequality follows by splitting at $x=n^{-1/q}$.
Therefore
\[
 \E\log L
 \le\frac{n^{1/q}}n\sum_{p,a}k_{p,a}\log p
 =\frac{n^{1/q}}n\sum_i\log c_i.
\]
At least one tuple of permutations has logarithm no larger than its
expectation, proving the result.
\end{proof}

\begin{lemma}[A positive coin bound]\label{coin}
Let $a_1,\ldots,a_q\le H$ be positive integers of greatest common
divisor $g$. Every multiple $C$ of $g$ with $C\ge qH+H^2$ has a
representation $C=\sum_{t=1}^q b_ta_t$ with all $b_t$ positive integers.
\end{lemma}
\begin{proof}
Subtract one copy of every generator. It remains to represent a
multiple $C'$ of $g$ with $C'\ge H^2$ by nonnegative coefficients.
Divide the generators by $g$, and let $a$ be their minimum.
The directed residue graph modulo $a$, with each generator allowed
as a step, is strongly connected because the scaled gcd is one.
A shortest path from zero to any residue has at most $a-1$ edges.
It therefore represents that residue by a nonnegative combination
of scaled generators of total value at most $(a-1)H/g$.
In the original units this value is at most $(a-1)H\le H^2$.
Choose the residue of $C'/g$ and add copies of the minimum generator
to reach $C'$. Restoring the initially subtracted generators makes
all coefficients positive.
\end{proof}

\begin{theorem}\label{balance}
Let $s$ be a permutation-fair word on $n\ge2$ letters, with
multiplicities $c_i$, length $A=\sum_i c_i$, and $H=\max_i c_i$.
There are $q=\lceil\log_2 n\rceil$ permutation-fair words
$s_1,\ldots,s_q$, each obtained by relabeling $s$ and independently
blowing up occurrences of each letter, whose combined multiplicity
of every letter is the same integer $C$, with
\[
 C\le qH+2H^2,
 \qquad \sum_t|s_t|=nC\le3A^3.
\]
\end{theorem}
\begin{proof}
Use Lemma~\ref{gcd}. Since $n^{1/q}\le2$ and the geometric mean is
at most $H$, its resulting $L$ satisfies $L\le H^2$.
Let $C$ be the least multiple of $L$ no smaller than $qH+H^2$.
Then $C\le qH+2H^2$.
For each coordinate $i$, its $q$ generators $c_{\sigma_t(i)}$
have gcd dividing $L$, hence dividing $C$. Lemma~\ref{coin} gives
positive integers $b_{ti}$ with
\[
 \sum_{t=1}^q b_{ti}c_{\sigma_t(i)}=C.
\]
Relabel $s$ according to $\sigma_t$, then replace every occurrence
of letter $i$ by $b_{ti}$ consecutive copies. Every full-pattern
count is multiplied by $\prod_i b_{ti}$, so each resulting word
$s_t$ is still permutation-fair. Its $i$-multiplicity is
$b_{ti}c_{\sigma_t(i)}$, giving the asserted common total.
Finally $n,q,H\le A$, so $n(qH+2H^2)\le3A^3$.
\end{proof}

\paragraph{The remaining geometric obstruction.}
The concatenation $s_1\cdots s_q$ is balanced in multiplicities,
but its order probabilities include selections spanning different
blocks. Their contribution is not controlled by the arithmetic
identities above. A construction that eliminated these cross-block
defects at cost exponential only in the number $q$ of profiles would
make the lemma relevant to polynomial equalization. No such
construction is established here. There is an additional concrete
barrier: if only these $q$ profile directions are reordered or repeated,
the block-profile matrix has rank at most $q$. The independently proved
continuous CDF-rank theorem in
\path{../size_bounds/two_dimensional_equalization_barrier.tex} then
forces at most $q$ distinct normalized letter profiles, with at most
one profile allowed to repeat. Thus at least $n-q+1$ letters must
share one profile whenever such a concatenation is three-wise fair.
For generic output profiles this fails, regardless of the number of
repetitions. Any eventual geometric use of the arithmetic lemma would
therefore need new profile directions or additional shared structure.
\end{document}
